% Clase 10 --- Capacitor de placas cilíndricas (§3).
% Dos vistas del mismo objeto: el corte transversal muestra la simetría que
% justifica Gauss; la vista lateral muestra las dos cosas que hacen falta para
% cerrar la cuenta y que en el corte no se ven: que las tapas no aportan flujo
% y que la gaussiana de altura H sólo encierra la fracción H/L de la carga.
\begin{center}
\begin{tikzpicture}[scale=0.92]

  % =============== (a) corte transversal ===============
  \begin{scope}
    \draw[figblue, line width=1.6pt] (0,0) circle (2.0);
    \foreach \a in {15,60,...,345} {\node[figlbl, figblue, scale=0.7] at (\a:1.78) {$-$};}
    \draw[figred, line width=1.4pt, fill=figred!8] (0,0) circle (0.6);
    \foreach \a in {30,105,...,340} {\node[figlbl, figred, scale=0.7] at (\a:0.8) {$+$};}

    \draw[figaux, line width=0.9pt] (0,0) circle (1.35);
    \node[figlblaux] at (128:1.55) {$S$};

    \foreach \a in {215,260,305,350} {\draw[figvecres] (\a:0.66) -- (\a:1.9);}
    \node[figlbl, figred, anchor=north] at (0,-2.35) {$\vec E$ radial};

    \draw[figvecaux] (0,0) -- (175:0.6);
    \node[figlblaux, anchor=south] at (175:0.33) {$A$};
    \draw[figvecaux] (0,0) -- (60:1.35);
    \node[figlblaux, anchor=north west] at (60:1.05) {$r$};
    \draw[figvecaux] (0,0) -- (40:2.0);
    \node[figlblaux, anchor=north west] at (40:1.55) {$B$};

    \node[figlbl, align=center] at (0,-3.5) {(a) corte transversal};
  \end{scope}

  % =============== (b) vista lateral ===============
  \begin{scope}[xshift=9.6cm]
    % cilindro externo
    \draw[figblue, line width=1.4pt] (-1.5,-2.3) -- (-1.5,2.3);
    \draw[figblue, line width=1.4pt] (1.5,-2.3) -- (1.5,2.3);
    \draw[figblue, line width=1.0pt, dash pattern=on 3pt off 2pt] (0,2.3) ellipse (1.5 and 0.34);
    \draw[figblue, line width=1.0pt, dash pattern=on 3pt off 2pt] (0,-2.3) ellipse (1.5 and 0.34);
    \node[figlbl, figblue, anchor=east] at (-1.6,1.75) {$-Q$};

    % cilindro interno: la parte encerrada por S, resaltada
    \fill[figred!22] (-0.45,-0.7) rectangle (0.45,0.7);
    \draw[figred, line width=1.2pt] (-0.45,-2.3) -- (-0.45,2.3);
    \draw[figred, line width=1.2pt] (0.45,-2.3) -- (0.45,2.3);
    \draw[figred, line width=1.0pt] (0,2.3) ellipse (0.45 and 0.11);
    \node[figlbl, figred, anchor=west] at (0.6,1.75) {$+Q$};

    % superficie gaussiana: cilindro de radio r y altura H
    \draw[figgray, dashed, line width=0.8pt] (-1.05,-0.7) rectangle (1.05,0.7);
    \draw[figgray, dashed, line width=0.8pt] (0,0.7) ellipse (1.05 and 0.24);
    \draw[figgray, dashed, line width=0.8pt] (0,-0.7) ellipse (1.05 and 0.24);

    % flujo por la pared lateral
    \foreach \y in {-0.35,0.35} {
      \draw[figvecres] (0.5,\y) -- (1.0,\y);
      \draw[figvecres] (-0.5,\y) -- (-1.0,\y);
    }
    % en las tapas E es paralelo a la tapa: no hay flujo
    \draw[figvecres] (0.2,0.92) -- (0.85,0.92);
    \draw[figgray, line width=0.4pt] (0.9,0.94) -- (1.62,1.05);
    \node[figlbl, figred, anchor=west, align=left] at (1.65,1.05)
      {$\vec E$ \emph{en el plano}\\[-0.15em] de la tapa:\\[-0.15em]
       \footnotesize flujo nulo};

    % cotas H y L
    \draw[figvecaux] (-1.28,0) -- (-1.28,0.7);
    \draw[figvecaux] (-1.28,0) -- (-1.28,-0.7);
    \node[figlblaux, anchor=east] at (-1.33,0) {$H$};
    \draw[figvecaux] (4.15,0.25) -- (4.15,2.3);
    \draw[figvecaux] (4.15,-0.25) -- (4.15,-2.3);
    \node[figlblaux, anchor=west] at (4.2,0) {$L$};

    % qué carga queda encerrada
    \draw[figgray, line width=0.4pt] (-0.5,-0.35) -- (-1.68,-0.9);
    \node[figlbl, figred, anchor=east, align=right] at (-1.72,-0.95)
      {\footnotesize sólo esta fracción\\[-0.2em] \footnotesize queda dentro de $S$};

    \node[figlbl, align=center] at (1.2,-3.5) {(b) vista lateral};
  \end{scope}

\end{tikzpicture}\\[0.5em]
{\small\itshape La hipótesis $L \gg B$ es la que permite tratar el campo como
puramente radial y uniforme a lo largo del eje: sólo así el corte (a) representa
igual de bien cualquier altura del cilindro. Con eso, las tapas de la gaussiana
no aportan flujo ---$\vec E$ está contenido en ellas--- y todo el flujo sale por
la pared lateral, de área $2\pi r H$. La carga encerrada no es $Q$ sino la
fracción de altura $H/L$, y esa $H$ se cancela contra la del área: el resultado
es exactamente el de una \textbf{línea infinita} de densidad $\lambda = Q/L$
(Clase 6), como tiene que ser. Cerca de las tapas reales del capacitor esto deja
de valer.}
\end{center}
